%% kapa: exam=1 round=36 year=2025 period=1 form=A subject=civil_law subject_ko=민법 seq=4 number=4
%% source: Q-Net 감정평가사 기출문제 (cst00309) artlSeq=5231525 file=raw/1008/5231525/2025년 제36회 감정평가사 1차시험 1교시 문제지 원본.pdf
%% conversion: kordoc 4.24.51 -> tools/gen1.py (auto); review=proofread-sonnet; flags=none
\begin{question}{4}
부재자의 재산관리에 관한 설명으로 옳지 않은 것은? (다툼이 있으면 판례에 따름)
\begin{choices}
  \item 부재자로부터 위임받은 재산처분권이 있는 재산관리인은 그 재산을 처분함에 있어 법원의 허가를 받아야 한다.
  \item 부재자 소유의 부동산에 관한 재산관리인의 매매행위에 대한 법원의 허가결정은 기왕의 매매를 추인하는 방법으로도 할 수 있다.
  \item 재산관리인을 선임한 부재자의 생사가 분명하지 아니한 때에 법원은 이해관계인의 청구에 의하여 재산관리인을 개임할 수 있다.
  \item 법원이 선임한 재산관리인은 법원의 허가없이 보존행위로서 부동산 소유권이전등기 말소등기절차 이행청구를 할 수 있다.
  \item 법원이 선임한 재산관리인이 부재자의 재산을 처분하기 위해서는 특별한 사정이 없는 한 법원의 허가를 받아야 한다.
\end{choices}
\end{question}
\answer{1}
